\scanpage{18}
\uline{$\overline{\mathrm{XVIII}}$} -- 1-2.

\margin{OK.} l. 7 \qquad $K_0$ \unsure{entre|autre|attention} \unsure{les|lui} \unsure{mains|mais} $G$
\[
K(K_0,\varphi)=p^*(K_0)+p^*(\varphi)[\unsure{O|\mathcal O}(\unsure{1|\xi|{-}1})]
\]

\par\noindent\rule{.8\linewidth}{.2pt}

\margin{OK} ligne 1 du bas

\unsure{peutêtre|peut être|peut-être} \unsure{métre|mettre} $R^1p_*p^*G$

\par\noindent\hspace*{2em}\rule{.35\linewidth}{.2pt}

\margin{OK} 1.2.2.1 : j'ai mis du \unsure{tem|temps} à comprendre
\unsure{d'abr.|d'abord} \unsure{lire} il \unsure{devien|devient|dernier} \unsure{cd.|c.d.}.

\par\noindent\rule{.9\linewidth}{.2pt}

on \unsure{pourrait|pourrai} faire remarque
\[
\Gamma(S,R^1p_*G)=\frac{H^1(\unsure{P|\mathbb P}(\unsure{E|\mathscr E}),G)}{H^1(S,G)}
\]
\unsure{quand|que|qu'} il y a une section
% archive sheet number: 17
